\documentclass[10pt,a4paper]{amsart}
\usepackage{amsmath,amssymb,amsthm,graphicx}
% Retained minimal preamble for the Stacks Project chapter sites-cohomology.tex (Cohomology on Sites).
% Verbatim excerpt of the chapter preamble: Xy-pic setup (native preamble lines 14-18) and the native
% theorem declarations with the mathematical macros (native preamble lines 154-217). Copied mechanically;
% nothing added, removed or redefined. The wrapper class and the amsmath/amssymb/amsthm/graphicx base
% packages are supplied by the assembler, so the class/externaldocument part of the native preamble is omitted.
% For commutative diagrams we use Xy-pic
\usepackage[all]{xy}

% We use 2cell for 2-commutative diagrams.
\xyoption{2cell}
% Theorem environments.
%
\theoremstyle{plain}
\newtheorem{theorem}[subsection]{Theorem}
\newtheorem{proposition}[subsection]{Proposition}
\newtheorem{lemma}[subsection]{Lemma}

\theoremstyle{definition}
\newtheorem{definition}[subsection]{Definition}
\newtheorem{example}[subsection]{Example}
\newtheorem{exercise}[subsection]{Exercise}
\newtheorem{situation}[subsection]{Situation}

\theoremstyle{remark}
\newtheorem{remark}[subsection]{Remark}
\newtheorem{remarks}[subsection]{Remarks}

\numberwithin{equation}{subsection}

% Macros
%
\def\lim{\mathop{\mathrm{lim}}\nolimits}
\def\colim{\mathop{\mathrm{colim}}\nolimits}
\def\Spec{\mathop{\mathrm{Spec}}}
\def\Hom{\mathop{\mathrm{Hom}}\nolimits}
\def\Ext{\mathop{\mathrm{Ext}}\nolimits}
\def\SheafHom{\mathop{\mathcal{H}\!\mathit{om}}\nolimits}
\def\SheafExt{\mathop{\mathcal{E}\!\mathit{xt}}\nolimits}
\def\Sch{\mathit{Sch}}
\def\Mor{\mathop{\mathrm{Mor}}\nolimits}
\def\Ob{\mathop{\mathrm{Ob}}\nolimits}
\def\Sh{\mathop{\mathit{Sh}}\nolimits}
\def\NL{\mathop{N\!L}\nolimits}
\def\CH{\mathop{\mathrm{CH}}\nolimits}
\def\proetale{{pro\text{-}\acute{e}tale}}
\def\etale{{\acute{e}tale}}
\def\QCoh{\mathit{QCoh}}
\def\Ker{\mathop{\mathrm{Ker}}}
\def\Im{\mathop{\mathrm{Im}}}
\def\Coker{\mathop{\mathrm{Coker}}}
\def\Coim{\mathop{\mathrm{Coim}}}

% Boxtimes
%
\DeclareMathSymbol{\boxtimes}{\mathbin}{AMSa}{"02}

%
% Macros for moduli stacks/spaces
%
\def\QCohstack{\mathcal{QC}\!\mathit{oh}}
\def\Cohstack{\mathcal{C}\!\mathit{oh}}
\def\Spacesstack{\mathcal{S}\!\mathit{paces}}
\def\Quotfunctor{\mathrm{Quot}}
\def\Hilbfunctor{\mathrm{Hilb}}
\def\Curvesstack{\mathcal{C}\!\mathit{urves}}
\def\Polarizedstack{\mathcal{P}\!\mathit{olarized}}
\def\Complexesstack{\mathcal{C}\!\mathit{omplexes}}
% \Pic is the operator that assigns to X its picard group, usage \Pic(X)
% \Picardstack_{X/B} denotes the Picard stack of X over B
% \Picardfunctor_{X/B} denotes the Picard functor of X over B
\def\Pic{\mathop{\mathrm{Pic}}\nolimits}
\def\Picardstack{\mathcal{P}\!\mathit{ic}}
\def\Picardfunctor{\mathrm{Pic}}
\def\Deformationcategory{\mathcal{D}\!\mathit{ef}}

% --- appended 2026-10-12 (Chief): the project own definitions of the reference/slogan/history
% display environments, copied mechanically from the shared Stacks preamble lines 8-12
% (\usepackage{verbatim} plus the three \newenvironment lines, which the chapter preamble keeps
% verbatim from the project preamble). Nothing added, removed or redefined: these environments
% carry no mathematical content and the project itself typesets them as comments. Without these
% three lines a counted statement whose own text uses them fails with Environment slogan undefined.
\usepackage{verbatim}
\newenvironment{reference}{\comment}{\endcomment}
\newenvironment{slogan}{\comment}{\endcomment}
\newenvironment{history}{\comment}{\endcomment}
\usepackage{cleveref}
\crefname{originalreference}{original reference}{original references}
\Crefname{originalreference}{Original reference}{Original references}
\makeatletter
\newlabel{lemma-build-from}{{0GYX}{source}{}{sourcepointer-1}{}}
\newlabel{lemma-build-from@cref}{{[originalreference][0][]0GYX}{[1][0][]source}{}{}{}}
\makeatother
\title{Every object of the derived category has a representing subcomplex of bounded size inside any given representative}
\author{The Stacks Project Authors}
\date{stacks-project:sites-cohomology.tex:local-snapshot-20261011; GNU FDL 1.2 or later}
\begin{document}
\maketitle
\noindent\textit{Editorial note.} All delivered mathematical text is the author's own, copied from the pinned Stacks Project chapter file sites-cohomology.tex; no statement or proof step was rewritten, reordered or removed.\par\medskip
\noindent\textit{Editorial note.} The counted proof is the author's complete argument: the choice of bounded subsets representing the cohomology classes, the construction of the initial bounded subcomplex via the building lemma, and the inductive chain of subcomplexes with stabilizing kernels whose union solves the problem.\par\medskip
\noindent\textit{Editorial note.} The label lemma-qis-small names the counted statement, and the proof environment immediately following it in the source is delivered verbatim as its proof.\par\medskip
\noindent\textit{Editorial note.} The counted statement is carried by the author's own lemma environment in the chapter source.\par\medskip
\noindent\textit{Editorial note.} Pre-open gate: python operations/agent-authoring/tools/precard\_gate.py --source sites-cohomology.tex --label lemma-qis-small --start 11237 --end 11293 --local-id 3554 returned RESULT=OK (L3554 free, occupancy FREE, public interval clear). Candidate taken from lane-s4/halfB-priority.json.\par\medskip
\noindent\textit{Editorial note.} The statement and proof are copied from the pinned chapter file; the cited inputs are kept as external pointers with their canonical Stacks tags.\par\medskip
\noindent\textit{Allowed standard prerequisites (editorial).}\begin{itemize}
\item Building a subcomplex of bounded size containing prescribed subsets. Same chapter, the lemma constructing, from a cardinal kappa, a subcomplex of a given complex of O-modules of size at most max(kappa, -) containing prescribed subsets of the terms. Supplies, at each step of the induction, the subcomplex of bounded size containing the chosen subsets of the groups of sections.
\end{itemize}\medskip
\ifdefined\hypertarget\hypertarget{sourcepointer-1}{}\fi
\noindent\textit{Source reference.} Original reference 0GYX: sources/stacks-sites-cohomology/pinned-original-source.tex, source lines 11217--11217. Prerequisite pointer invoked at proof line(s) 11255; it names a standard fact of the same chapter that the counted argument uses, and every step of the counted proof built on it is retained verbatim in the delivered body.\par\medskip
\begin{lemma}
\label{lemma-qis-small}
In the situation above, there exists a cardinal $\kappa$ with the
following property: given a complex $\mathcal{F}^\bullet$ of
$\mathcal{O}$-modules representing an object $K$ of $D(\mathcal{O})$
there exists a subcomplex $\mathcal{H}^\bullet \subset \mathcal{F}^\bullet$
such that $\mathcal{H}^\bullet$ represents $K$ and such that
$|\mathcal{H}^\bullet| \leq \max(\kappa, |K|)$.
\end{lemma}

\begin{proof}
First, for every $i$ and $U$ we choose a subset
$\Omega^i_U \subset \Ker(\text{d} :
\mathcal{F}^i(U) \to \mathcal{F}^{i + 1}(U))$
mapping bijectively onto
$H^i(K)(U) = H^i(\mathcal{F}^\bullet(U))$.
Hence $|\Omega^i_U| \leq |K|$ as we may represent $K$ by
a complex whose size is $|K|$.
Applying Lemma \ref{lemma-build-from}
we find a subcomplex $\mathcal{S}^\bullet \subset \mathcal{F}^\bullet$
of size at most $\max(\kappa, |K|)$ containing $\Omega^i_U$
and hence such that $H^i(\mathcal{S}^\bullet) \to H^i(\mathcal{F}^\bullet)$
is a surjection of sheaves.

\medskip\noindent
We are going to inductively construct subcomplexes
$$
\mathcal{S}^\bullet =
\mathcal{S}_0^\bullet \subset
\mathcal{S}_1^\bullet \subset
\mathcal{S}_2^\bullet \subset \ldots \subset \mathcal{F}^\bullet
$$
of size $\leq \max(\kappa, |K|)$ such that
the kernel of $H^i(\mathcal{S}_n^\bullet) \to H^i(\mathcal{F}^\bullet)$
is the same as the kernel of
$H^i(\mathcal{S}_n^\bullet) \to H^i(\mathcal{S}_{n + 1}^\bullet)$.
Once this is done we can take
$\mathcal{H}^\bullet = \bigcup \mathcal{S}_n^\bullet$ as our solution.

\medskip\noindent
Construction of $\mathcal{S}_{n + 1}^\bullet$ given
$\mathcal{S}_n^\bullet$. For ever $U$ and $i$ let
$\Omega^{i - 1}_U \subset \mathcal{F}^{i - 1}(U)$ be a subset
such that $\text{d} : \mathcal{F}^{i - 1}(U) \to \mathcal{F}^i(U)$
maps $\Omega^{i - 1}_U$ bijectively onto
$$
\mathcal{S}_n^i(U) \cap
\text{Im}(\text{d} : \mathcal{F}^{i - 1}(U)\to \mathcal{F}^i(U))
$$
Observe that $|\Omega^i_U| \leq |K|$ because $\mathcal{S}_n^i(U)$
is so bounded. Then we get $\mathcal{S}_{n + 1}^\bullet$ by an
application of Lemma \ref{lemma-build-from} to the subsets
$$
\mathcal{S}^i(U) \cup \Omega^i_U \subset \mathcal{F}^i(U)
$$
and everything is clear.
\end{proof}

\end{document}
